\section*{The Propers: Detailed Commentary}

\subsection*{The Table that Tests Banquet Speech\enspace\properrefs{Gosp.}}

Luke places the parable at a Pharisee's table after a guest blesses the one who will eat bread in God's kingdom. The saying is therefore not distant banquet imagery. It tests the present hearer: admiration for the kingdom can coexist with refusal when the actual summons interrupts possession, investigation, or pleasure.

\sectionguard
\subsection*{The Collect: Fear and Love under Continuing Governance\enspace\properrefs{Coll.}}

The Collect does not ask the worshipper to choose between fear and love. It asks God to make both dispositions endure together toward his holy Name, and it grounds that endurance in God's continuing governance and prior establishment of the servants in his love.

\begin{threestable}{Latin movement}{What the prayer directly asks}{Precise sense}
\latin{timorem pariter et amorem} & Fear and love are held together rather than assigned to rival stages of Christian life. & The oration does not define either disposition as a passing emotion or reduce fear to panic.\\
\latin{fac nos habere perpetuum} & The worshippers ask God to make their response enduring. & Stability is petitioned as gift; it is not self-certified before the invitation is tested.\\
\latin{numquam tua gubernatione destituis} & God's never-failing governance supplies the reason for confidence. & Perseverance rests first on divine action without making the servants' answer irrelevant.\\
\latin{in soliditate tuae dilectionis instituis} & God establishes the governed within the solidity of his own love. & The prayer names the ground of stability; it does not claim that every present attachment is already rightly ordered.\\
\end{threestable}

The Collect therefore contributes more than a general request for good feeling. Before First John tests charity in material action and the Gospel tests desire at the moment of invitation, the Church asks that reverence and love be made durable by the God who does not abandon what he establishes.

\sectionguard
\subsection*{The Host and Supper: Four Patristic Readings\enspace\properrefs{Gosp.}}

The four witnesses converge on the feast as God's saving abundance disclosed in Christ, but they do not offer one role-by-role diagram. Tertullian, Augustine, and Cyril place the invitations within salvation history in different ways; in §§1--5 Gregory turns instead toward desire, conduct, and the urgency of the hearer's present answer.

\begin{fourstable}{Witness and locus}{Host, supper, and servant}{Historical movement}{Distinctive contribution}
Tertullian, \emph{Against Marcion} IV.31 & The Creator is the host and eternal life the abundant supper. & The patriarchs and fathers receive the first invitation; prophetic admonition summons their descendants; the poor are welcomed in continuity with Isaiah. & His anti-Marcionite purpose establishes one Creator and one economy across the Testaments; he does not assign every role.\\
St. Augustine, \emph{Sermon 62}.1 & Christ the Mediator is host; the supper is his saving gift, explicitly including the Eucharistic food offered daily; the apostles issue the summons. & Prophets first invited Israel; Christ's death prepares the feast; apostolic preaching announces that all is ready. & Augustine concentrates the whole economy in Christ and then turns quickly to the hearer's interior readiness.\\
St. Cyril of Alexandria, \emph{Commentary on Luke}, Sermon 104 & The Father prepares the supper in Christ; the one servant signifies the Son's incarnation and ministry. & Synagogue leaders refuse, ordinary Jewish hearers are gathered from the city, and the apostolic mission reaches the Gentile roads. & Cyril names forgiveness, cleansing, the Spirit, adoption, and the kingdom as the feast's gifts. His servant mapping differs from Augustine's without denying apostolic agency.\\
St. Gregory the Great, \emph{Homiliae in Evangelia} II.36 & The great supper is the eternal sweetness offered by God. & In §§1--5 he turns from narrative roles to the psychology of desire and the contradiction between faith and conduct. & Spiritual goods may initially seem unattractive but deepen desire when tasted; bodily pleasures often reverse that pattern. In §§1--5 Gregory reads the supper morally.\\
\end{fourstable}

St.~Anthony of Padua's medieval sermon on the same Sunday refuses to isolate the banquet from the appointed Epistle and Introit. He calls Christ the man who prepares both the supper of repentance and the supper of glory; the Passion opens the kingdom, while brotherly love introduces the guest into eternal life. His treatment of the excuses is more ascetically expansive than Augustine's and often addresses preachers and religious directly. Its durable contribution is narrower: the invitation is prior grace, lawful goods become exclusionary when they govern desire, and charity toward the brother is already a mode of entering the feast. Anthony receives the Gospel, Epistle, and Introit together in his own medieval preaching of the Sunday.

\sectionguard
\subsection*{Three Excuses and the Love Beneath Them\enspace\properrefs{Ep., Gosp.}}

In \emph{Sermon 62}.2--7 Augustine joins the three excuses to the threefold worldly desire of 1 John 2:16. His allegory does not condemn land, sense knowledge, or marriage. It exposes the governing love that can turn a lawful good into a courteous refusal of the giver.

\begin{fourstable}{Excuse}{Augustine's interpretation}{Johannine disorder}{How this Sunday's Epistle tests it}
An estate bought and inspected & The appetite to dominate through possession: the field of rule displaces the kingdom one is invited to receive. & Pride of life. & One who closes worldly goods against a needy brother reveals that possession has not become charity.\\
Five yoke of oxen to be proved & The five bodily senses and a curiosity that grants authority only to what can be experimentally inspected. & Lust or curiosity of the eyes. & The passage from death to life is known through brotherly love, a spiritual reality that nonetheless becomes visible in deed.\\
A wife newly married & Fleshly delight made ultimate, so that an intimate created good becomes exemption from the Creator's summons. & Lust of the flesh. & Christ's laying down his life supplies the form of love: communion turns outward in costly self-gift rather than closing around pleasure.\\
\end{fourstable}

\begin{studybox}{The reach of Augustine's allegory}
\begin{itemize}
\item The literal guests are culpable for refusing an invitation already accepted; Augustine's symbolic identifications read them spiritually, and Luke's words keep that literal sense.
\item Property, empirical inquiry, bodily life, and marriage remain genuine goods. Disorder appears when any good becomes a reason the giver may no longer command.
\item The threefold parallel with 1 John 2:16 belongs to Augustine's sermon; linking this Sunday's separate Epistle from 1 John 3 to the parable is a reading of the Missal's pairing, not his.
\end{itemize}
\end{studybox}

Gregory the Great sharpens the diagnosis. One may confess faith and still contradict the banquet by conduct. Bodily pleasures can be intensely desired before possession and leave disgust afterward; spiritual delight is often resisted before it is tasted but enlarges desire as it is possessed. The problem is not simply a bad decision at one moment. The will has been trained to experience the lesser good as urgent and the supper as postponable.

\sectionguard
\subsection*{The Poor, the Roads, and ``Compel Them''\enspace\properrefs{Ep., Gosp.}}

The host's anger does not cancel generosity. It redirects the servants toward the poor, feeble, blind, and lame, then beyond the city into roads and hedges until the house is filled. Patristic interpretations of these groups illuminate the history of mission, and Cyril and Augustine read the command to compel differently.

\begin{threestable}{Narrative turn}{Cyril, Sermon 104}{Augustine, \emph{Sermon 62}.8}
City streets and lanes & People within Israel who lacked the prestige of the rejecting synagogue leaders are gathered into Christ's gifts. & Those who know their need replace guests confident in their sufficiency.\\
Roads and hedges & The command reaches the Gentiles through the universal apostolic mission; Cyril closes with Matthew 28. & Roads signify Gentile wandering; hedges signify divided or obstructed hearers whom Augustine associates with heresy.\\
\latin{Compelle intrare} & Divine urgency and drawing operate without making faith involuntary; Cyril insists on willing belief. & Augustine distinguishes outward pressure from the inward willing that must follow, in a passage later used in coercive controversies.\\
\end{threestable}

\begin{studybox}{\latin{Compelle intrare} and willing faith}
\latin{Compelle intrare} is no warrant for violence, forced sacramental profession, or confidence that exterior conformity equals faith. Cyril's insistence on voluntary belief and the parable's concern with genuine acceptance rule out that simplification. Augustine's account belongs to its late-antique polemical history and is not a timeless pastoral technique.
\end{studybox}

The widening invitation also rebukes contempt. The disabled and poor are not props in someone else's allegory; within Luke's narrative they are actual persons the host commands to be brought in. The Epistle's concrete demand to open one's goods to the brother in need tests every spiritual reading of the parable: a sound one strengthens that demand rather than evading it.

\sectionguard
\subsection*{First John: Sacrament, New Birth, and Charity in Deed\enspace\properrefs{Ep.}}

The Epistle does more than provide an ethical application after the parable. It distinguishes death and life, hatred and love, the world's goods and the brother's need, speech and deed, Christ's self-gift and ours. Augustine's \emph{Homily 5 on First John} follows that internal order.

\begin{fourstable}{Augustine's locus}{Direct exegesis of 1 John 3}{Sacramental precision}{Demand upon the invited guest}
Section 6 & Baptism marks the new person and remission of sins, but the heart must possess charity. & The sacramental mark is real; lack of charity does not make it imaginary, but it exposes a life not answering its grace. & Presence at the feast cannot be reduced to bodily access or verbal identity.\\
Section 10 & Passing from death to life is known through love of the brethren. Augustine compares winter's apparently dead tree whose living root remains hidden. & Charity can be real before its mature fruit is fully visible, because divine life is rooted within. & The world may not recognize the life it cannot see; the Christian must let the hidden root issue in action.\\
Sections 11--12 & Christ's laying down his life is charity's perfection; opening one's means to the needy is its serious beginning. & The same blood is the price of every redeemed brother, so sacramental communion forbids a merely private valuation of grace. & Love in deed starts with available material mercy and tends toward Christ-shaped self-gift.\\
\end{fourstable}

Anthony makes the Epistle's bodily test especially concrete. After naming the Introit and 1~John 3:13--18 alongside Luke 14, he treats love of the brethren as entrance into the banquet and reads the closed heart as a refusal of the Host who first stretched out his hand on the Cross. Augustine supplies the more exact sacramental distinction between baptismal mark and living charity; Anthony shows how medieval preaching carried that distinction into examination of possession, sense, power, and material mercy.

\begin{massbox}{Neither empty badge nor self-made worthiness}
Augustine neither dismisses baptism when charity is lacking nor treats the sacramental mark as a guarantee that conduct no longer matters. The invited person receives a gift he did not prepare and is made responsible to live from it. Charity is evidence of passage into life because the life itself comes first from God.
\end{massbox}

\sectionguard
\subsection*{The Psalms Retrain the Guest's Voice\enspace\properrefs{Int., Grad., All., Off., Comm.}}

Each chant verse has a substantial Augustinian exposition. Augustine interprets the appointed biblical words directly, though not their later arrangement in this Mass. The movement from rescue, through truthful speech and hope, to petition and thanksgiving is a reading of the chants in the order this Mass gives them.

\begin{fourstable}{Chant and verse}{Augustine's direct exegesis}{Conversion of the hearer}{Locus}
Introit, Psalm 17:19--20 & In \emph{Exposition on Psalm 17}, 20--21, the broad place is faith's spiritual breadth; deliverance rests on God's prior delight and mercy, with good will and clean deeds following the first gift. & The guest enters room created by grace before being sent to make room for another. & Psalm 18 in modern numbering; the missal follows the Vulgate.\\
Gradual, Psalm 119:1--2 & In \emph{Exposition on Psalm 119}, 2--8, the cry begins spiritual ascent. Deceitful tongues include counselors whose praise or advice draws the pilgrim away from Christ; God's word can become an arrow that converts an enemy. & The hearer asks to be freed from the speech that will soon appear as polished excuse, then learns to sing the psalm with the heart and life. & Psalm 120 in modern numbering.\\
Alleluia, Psalm 7:2 & In \emph{Exposition on Psalm 7}, 2, hope is placed in God against persecutors and, beneath them, the enemy figured as a tearing lion. & The world's hatred named by the Epistle is answered by trust and patient refusal to repay evil, not by self-rescue. & The devil-image belongs to Augustine's psalm exegesis, not to Luke's banquet characters.\\
Offertory, Psalm 6:5 & In \emph{Exposition on Psalm 6}, 5, ``Turn, O Lord'' can be heard as the plea that God make the sinner turn and assist a difficult return; salvation is sought for mercy, not merit. & The excuse-maker ceases explaining himself and asks God for the conversion he cannot perfect unaided. & Augustine expounds the psalm's imperative, not the Secret.\\
Communion, Psalm 12:6 & In \emph{Exposition on Psalm 12}, 6, hope rests in divine mercy rather than self-stability; song belongs to the heart and is chanted publicly in works. & The tongue rescued from deceit answers received benefits with gratitude embodied in deed and truth. & Psalm 13 in modern numbering.\\
\end{fourstable}

\sectionguard
\subsection*{The Secret: Purified for the Conduct of Heaven\enspace\properrefs{Sec., Postcomm.}}

The Secret gives the moral demand a sacramental grammar. Its subject is the oblation; its repeated object is the worshippers. The prayer asks neither that they purify themselves nor merely that their existing conduct be approved.

\begin{threestable}{Latin hinge}{Syntactical force}{Place in the Sunday's movement}
\latin{Oblatio ... tuo nomini dicanda} & The oblation is to be dedicated to the Lord's Name. & The Collect's fear and love of the holy Name reach the altar as consecrated action.\\
\latin{nos ... purificet} & The oblation is asked to purify \emph{us}. & Guests do not arrive as self-certified worthy recipients; divine gift cleanses those it summons.\\
\latin{et de die in diem ... transferat} & The same \latin{nos} naturally remains the object: may the oblation transfer us day by day. & Transformation is repeated and progressive, not exhausted by one moment of enthusiasm.\\
\latin{ad caelestis vitae ... actionem} & The goal is the action, conduct, or way of life proper to heavenly life. & The banquet's anagogical promise enters present behavior; heaven begins to govern how the guest acts toward the brother.\\
\end{threestable}

\begin{sourcecard}{St. Augustine --- \emph{City of God} X.5--6, 20}
Augustine does not comment on this Secret, but his sacrificial account clarifies its logic. Every work ordered to holy union with God is true sacrifice; mercy toward the neighbor belongs to that order; body and soul are offered as the mind is transformed; and the whole redeemed city offers itself through Christ the High Priest. The altar's sacrament teaches the Church to become, through him, the self-offering it celebrates.
\end{sourcecard}

The relation to the Epistle is exact without being mechanical. First John says the faithful have passed from death to life and gives brotherly love as evidence. The Secret asks that sacramental oblation continue transferring them, day by day, into heaven's mode of action. An accomplished passage and an unfinished formation coexist because baptismal life has been given as a principle that must mature in charity.

The Postcommunion completes the same order. Having received the sacred gifts, the Church asks that repeated participation in the mystery increase salvation's effect. Frequency is not treated as bare repetition; it is ordered toward growth. The supper that was ready before the guest arrived becomes nourishment whose effect must appear in a life opened to God and neighbor.
